% !PRINTED 264
\weberpage{264}
From this the degrees of the elements of our group can easily be
determined. We have three cases to distinguish here.

I.~\gesperrt{$a + d = \pm 2$.} \gesperrt{Every substitution of this kind can be
transformed into the normal form $T$. But for every
exponent $r$}
\[
T^r = \begin{pmatrix} 1, & r\,t \\ 0, & 1 \end{pmatrix},
\]
\gesperrt{and consequently the degree of this substitution is $p$, except
when $t = 0$, that is, when $T$ is the identical substitution.}

The number of elements of this kind is easy to determine.
If we first take $c = 0$, we can choose $b$ arbitrarily,
and for $a, d$ we obtain from $a + d = \pm 2$, $a\,d = 1$ the two
determinations $a = d = \pm 1$. This gives $2p$ forms, the identical
substitution included. If we next take $a$ and $c$
arbitrary, but $c$ different from zero, which gives $p\,(p-1)$ possibilities,
then from $a\,d - b\,c = 1$ it follows that
\[
b = -c^{-1} + a\,d\,c^{-1},
\]
and for each $a$, $d$ can be determined in two ways. The
total number $2\,p^2$ that we obtain in this way must still be halved, because
here the substitutions occur as two different ones in which
all signs are opposite; and we thus obtain $p^2$ substitutions
of this kind (the identical one included).

II.~\gesperrt{$(a+d)^2 - 4$ a quadratic residue of $p$. Such a
substitution can be transformed into the normal form $S$
with real $\sigma$.}

If we understand by $\gamma$ a primitive root of the field
$\mathfrak{E}_{2,p}$, then
\[
\gamma^{\frac{p^2-1}{2}} = -1,
\]
and we can determine $r$ so that
\[
\sigma = \gamma^{r(p+1)}, \quad \sigma^{\frac{p-1}{2}} = (-1)^r
\]
holds (\S.\,67). By (5) we then obtain
\weq{9}{a + d = \gamma^{r(p+1)} + \gamma^{-r(p+1)},}
and since by (5) the two values $\sigma, \sigma^{-1}$ are, apart from
their order, determined by $a+d$, two values
of the expression (9) become equal or opposite to one another
only when the corresponding values, with positive or
%CONT>>
